% When submitting your files, remember to upload this *tex file, the pdf generated with it, the *bib file (if bibliography is not within the *tex) and all the figures.

% NB logo1.jpg is required in the path in order to correctly compile front page header %%%

\documentclass[utf8]{FrontiersinHarvard} % for articles in journals using the Harvard Referencing Style (Author-Date), for Frontiers Reference Styles by Journal—https://zendesk.frontiersin.org/hc/en-us/articles/360017860337-Frontiers-Reference-Styles-by-Journal
\usepackage{url,hyperref,lineno,microtype,subcaption}
%\usepackage[]{natbib, booktabs, graphicx, tabularx, makecell}
\usepackage{booktabs}
\usepackage{lscape, float, multirow, adjustbox}
\usepackage[T1]{fontenc}
\usepackage[onehalfspacing]{setspace}

\linenumbers


% Leave a blank line between paragraphs instead of using \\


\def\keyFont{\fontsize{8}{11}\helveticabold }
\def\firstAuthorLast{Morris {et~al.}} %use et al only if is more than 1 author
\def\Authors{Joanna Morris\,$^{1,2,*}$, Emma Kealey\,$^{2,3}$, Frankie Greene\,$^{1,4}$, Joemari Pulido\,$^{2,5}$, Tess Rooney\,$^{2,6}$, Ethan Moore\, $^{2}$, Crismar Ramos-Marte\, $^{2}$, Natalia Alzate\,$^{2}$}
% Affiliations should be keyed to the author's name with superscript numbers and be listed as follows—Laboratory, Institute, Department, Organization, City, State abbreviation (USA, Canada, Australia), and Country (without detailed address information such as city zip codes or street names).
% If one of the authors has a change of address, list the new address below the correspondence details using a superscript symbol and use the same symbol to indicate the author in the author list.
\def\Address{$^{1}$School of Cognitive Science, Hampshire College, Amherst, MA, USA \\
	$^{2}$Department of Psychology, Providence College, Providence, RI, USA \\
	$^{3}$Department of Linguistics, University of Southern California, Los Angeles, CA, USA \\
	$^{4}$Oregon Health \& Science University and Portland State University School of Public Health, Portland, OR, USA \\
	$^{5}$Department of Cognitive, Linguistic \& Psychological Sciences, Brown University, Providence, RI, USA \\
	$^{6}$Columbia University School of Nursing, New York, NY, USA}
% The Corresponding Author should be marked with an asterisk
% Provide the exact contact address (this time including street name and city zip code) and email of the corresponding author
\def\corrAuthor{Joanna Morris}

\def\corrEmail{jmorris6@providence.edu}


\begin{document}
	\onecolumn
	\firstpage{1}
	
	\title[Running Title]{Reading skill and the time course of morphological processing: electrophysiological evidence for stage-specific modulation} 
	
	\author[\firstAuthorLast ]{\Authors} %This field will be automatically populated
	\address{} %This field will be automatically populated
	\correspondance{} %This field will be automatically populated
	
	\extraAuth{}% If there are more than 1 corresponding author, comment this line and uncomment the next one.
	%\extraAuth{corresponding Author2 \\ Laboratory X2, Institute X2, Department X2, Organization X2, Street X2, City X2 , State XX2 (only USA, Canada and Australia), Zip Code2, X2 Country X2, email2@uni2.edu}
	
	
	\maketitle
	
	
	
	
\begin{abstract}
		
		%%% Leave the Abstract empty if your article does not require one, please see the Summary Table for full details.
	\section{}
	Skilled reading of morphologically complex words depends on both high-quality lexical representations and sensitivity to morphological structure, yet how individual differences in these capacities shape the time course of morphological processing remains poorly understood. Event-related potentials (ERPs) and reaction times were recorded during a lexical decision task on words and morphologically complex nonwords varying in base frequency, morphological family size, and morphological complexity. Individual differences were quantified via principal components analysis of vocabulary knowledge, spelling ability, and print exposure, yielding two orthogonal dimensions: Language Proficiency and Orthographic Sensitivity. Behaviorally, response times were driven primarily by item-level morphological structure, with individual differences affecting overall processing speed but not the pattern of morphological effects. ERP measures revealed a strikingly different picture: Orthographic Sensitivity selectively modulated the N250, reflecting early form-based parsing, while Language Proficiency selectively modulated the N400, reflecting later semantic integration. This double dissociation across processing stages indicates that the distinction between form-based and meaning-based morphological processing—proposed by competing theoretical accounts on computational grounds—has a genuine temporal neural correlate that varies across readers. The findings demonstrate that individual differences in reading skill are expressed in the dynamics of neural processing rather than in overt response times, and underscore the importance of electrophysiological measures for understanding how morphological structure is parsed during skilled reading.
		
		
	\tiny
	\keyFont{ \section{Keywords:} morphological processing, individual differences, event-related potentials, N250, N400, lexical decision} %All article types—you may provide up to 8 keywords; at least 5 are mandatory.
\end{abstract}
	
	\section{Introduction}
	
	Reading is often described as "the process of gaining meaning from print" \citep[][p.~34]{raynerHowPsychologicalScience2001}. Although comprehension ultimately depends on sentence- and discourse-level integration, efficient access to individual word forms and meanings remains a foundational component of skilled reading. Crucially, readers vary considerably in this ability \citep{perfettiReadingAbilityLexical2007, castlesHowDoesOrthographic2006, stanovichExposurePrintOrthographic1989}. Understanding what drives this variation requires identifying variables that are selectively diagnostic of the mechanisms underlying word recognition. Two such variables, base frequency (also referred to as stem frequency in the literature) and morphological family size, have proven particularly central to this goal \citep{taftMorphologicalDecompositionReverse2004, taftRecognitionAffixedWords1979, dejongMorphologicalFamilySize2000, bertramEffectsFamilySize2000}. Base frequency, the cumulative token frequency of a base morpheme across all its surface forms, and morphological family size, the number of distinct words sharing that base, both facilitate visual word recognition, yet their effects are independent and have been argued to reflect different levels of morphological representation: base frequency has been argued to index form-based access, reflecting the cumulative ease with which a base morpheme's orthographic form is processed across its surface variants, while family size has been argued to index the richness of semantic representation, reflecting the breadth of meaning-level connections activated when a morphological base is encountered \citep{schreuderHowComplexSimplex1997}. The present study exploits this dissociation to ask how readers who differ in the precision of their orthographic and semantic knowledge draw differentially on form-based and meaning-based information during morphological processing.
	
	Properties of the English writing system make sensitivity to morphological structure particularly consequential for skilled reading. English exhibits irregular and context-sensitive grapheme–phoneme mappings, rendering phonological decoding alone an insufficient basis for fluent word recognition \citep{zieglerReadingAcquisitionDevelopmental2005, castlesHowDoesOrthographic2006}. At the same time, English preserves regularities at higher representational levels, particularly in its morphological structure; words that share a stem or affix often exhibit consistent orthographic patterns and systematic semantic relationships (e.g., \textit{heal–health}, \textit{teach–teacher}). Skilled reading in English therefore requires sensitivity not only to phonological correspondences but also to recurring morphological structure that links form and meaning, and readers differ in how deeply and precisely that structure is represented \citep{bryantMorphologySpellingWhat2005, arciuliReadingStatisticalLearning2018, treimanStatisticalLearningSpelling2018}.
	
	Characterising how readers differ in their sensitivity to base frequency and family size, however, first requires clarity about what these variables measure. This question has proven theoretically contentious because the answer depends on the architecture one ascribes to the word recognition system. Whether skilled readers decompose complex words into constituent morphemes, access whole-word representations directly, or exploit learned statistical regularities carries different implications for what each variable measures: whether it reflects the strength of a stored morpheme entry, the output of a parsing procedure, or an emergent property of a statistical learning system. Crucially, what each variable measures determines what kind of reader variation should modulate its effect. The Lexical Quality Hypothesis \citep{perfettiLexicalQualityHypothesis2002, perfettiReadingAbilityLexical2007, perfettiLexicalQualityRevisited2017} provides a principled framework for characterising that variation: lexical representations vary across readers in how completely and precisely a word's orthographic and semantic constituents are specified, and it is this variation in representational quality that determines how strongly each variable influences processing. Because base frequency indexes how precisely the orthographic form of the base morpheme is specified in the reader's lexicon, and family size indexes how richly and flexibly its meaning is represented across morphological contexts, readers who differ in orthographic precision versus semantic breadth should differ selectively in their sensitivity to each variable. If form-based and semantic processing are moreover temporally dissociable, those differences should emerge at distinct stages in the processing cascade.
	
	
	\subsection{Base Frequency, Morphological Family Size, and Competing Theoretical Accounts}
	
	Both base frequency and morphological family size facilitate visual word recognition, and their effects are independent when appropriately controlled \citep{fordDerivationalMorphologyBase2010}. The theoretical significance of this dissociation is that it constrains what any model of morphological processing must explain. Current models converge on the prediction that base frequency effects arise earlier in processing than family size effects, but differ considerably in how they explain this, and consequently in what they propose about the nature of the representations involved and the source of individual differences in sensitivity to each variable.
	
	Dual-route and multi-stage accounts agree that the two variables index distinct levels of processing, but disagree on the nature of those levels. In classical dual-route frameworks \citep{schreuderModelingMorphologicalProcessing1995, dejongMorphologicalFamilySize2000}, base frequency reflects the ease of accessing a stored morpheme entry via a decompositional route, while family size reflects the richness of the semantic network surrounding the base, attributed to the whole-form route. Early automatic decomposition accounts \citep{taft_lexical_1975, taftMorphologicalDecompositionReverse2004} reconceptualise this dissociation in temporal rather than architectural terms: decomposition is not one route among several but an obligatory initial parsing step that precedes semantic access. On the morpho-orthographic formulation of this view \citep{rastleMorphologicalDecompositionBased2008, longtinMorphologicalDecompositionEarly2005, morrisSemanticTransparencyMasked2007}, an intermediate level of representation mediates between sublexical letter clusters and whole-word forms, and base frequency reflects the strength of activation at this stage. Family size effects are then linked to later semantic processing, consistent with the finding that opaque relatives contribute less to the family size effect than transparent ones \citep{schreuderHowComplexSimplex1997, dejongMorphologicalFamilySize2000}. The discriminative learning framework \citep[NDL/LDL;][]{chuangDiscriminativeLearningLexicon} offers a more radical reinterpretation. On this account, lexical processing is modelled as learned form-to-meaning mappings driven by prediction error rather than by decompositional morpheme retrieval. Base frequency reflects the accumulated strength of the n-gram cues that spell out the base across all its surface forms, while family size reflects the type diversity of the training signal; a larger family causes shared cues to converge on a more robustly specified semantic core. Crucially, both effects arise from the same learning mechanism, rendering the dissociation between them a consequence of different statistical properties of the input rather than evidence for distinct representational levels or processing routes.
	
	These accounts converge on the prediction that base frequency and family size should exert independent effects, with base frequency effects arising earlier in processing than family size effects. Where they diverge is in the nature of the representations involved—symbolic morpheme entries, learned orthographic patterns, or weighted form-to-meaning mappings—and consequently in how they explain the source of individual differences in sensitivity to each variable. Dual-route and morpho-orthographic accounts imply that readers with more precise sublexical orthographic representations should show greater sensitivity to base frequency, since this variable indexes activation strength at the form-parsing stage, while readers with richer semantic networks should show greater sensitivity to family size, since this variable indexes the breadth of meaning-level connectivity. The discriminative learning framework leads to a similar expectation but for different reasons: individual differences in statistical learning history determine the specificity of both orthographic cue weights and semantic vectors, so readers with more extensive print exposure should show stronger and more differentiated sensitivity to both variables. Crucially, however, all three accounts are consistent with the view that form-level and meaning-level sensitivity should be dissociable, and that this dissociation should be traceable to distinct processing stages—early for base frequency, later for family size.
	
	
	\subsection{Individual Differences in Morphological Processing}
	
	Traditional models of word recognition have typically characterised processing at the level of the skilled reader in general, abstracting over individual differences in favour of a uniform cognitive architecture. Yet readers vary considerably in how they process written words, and understanding this variation offers a window into the representations and processes underlying skilled reading. The Lexical Quality Hypothesis, introduced above, provides a principled framework for understanding this variation. Critically, readers do not differ along a single dimension of lexical quality. Andrews and colleagues \citep{andrewsIndividualDifferencesSkilled2012, andrewsMorphologicalPrimingStronger2013} have demonstrated that individual differences in orthographic and semantic knowledge are dissociable and predict distinct patterns of morphological processing. Readers with a semantic profile, defined by relatively higher vocabulary than spelling ability, show robust priming for transparent morphological pairs but little priming for opaque or form-only pairs, particularly for slower responses. Readers with an orthographic profile, defined by relatively higher spelling than vocabulary, show sustained priming for opaque pairs that is at least as strong as for transparent pairs. \citet{andrewsMorphologicalPrimingStronger2013} interpret this pattern as evidence that a systematic component of individual differences among skilled readers is explicable in terms of the balance between orthographic and semantic processing; readers differ not merely in overall processing efficiency but in the relative weight they assign to form-based versus meaning-based information during lexical access.
	
	The present study builds directly on this finding. If base frequency and family size index form-based and meaning-based aspects of lexical representation respectively, then readers who differ in their balance of orthographic and semantic processing should differ predictably in their sensitivity to each variable. Moreover, if these differences have a temporal structure, with form-based processing preceding semantic processing in the recognition cascade, they should be traceable to distinct ERP components. To capture the full range of this variation continuously rather than as a categorical split, the present study derives two orthogonal individual difference dimensions, Language Proficiency and Orthographic Sensitivity, from a principal components analysis of vocabulary knowledge, spelling ability, and print exposure, following the approach of \citet{andrewsMorphologicalPrimingStronger2013}. If base frequency and family size are differentially modulated by these two dimensions, and if those modulations emerge at distinct stages in the ERP, this would constitute direct electrophysiological evidence that orthographic and semantic processing represent a genuine and temporally structured axis of variation in the human reading system.
	
	
	\subsection{Neural Indices of Morphological Processing}
	
	To track whether base frequency and family size effects emerge at distinct temporal stages, and whether individual differences in orthographic and semantic processing modulate those stages selectively, the present study uses two well-established ERP components as dependent measures—the N250, as an index of sublexical processing, and the N400, as an index of lexical-semantic access. These components have been argued to reflect sequential stages in the cascade from orthographic decoding to meaning retrieval during visual word recognition \citep{graingerWatchingWordGo2009}.
	
	The N250 is a negative-going deflection peaking at approximately 250 ms post-stimulus onset. \citet{holcombTimeCourseVisual2006} propose that it reflects the stage at which abstract sublexical orthographic representations—letters and letter clusters—are mapped onto higher-level form representations, prior to full lexical access. Evidence for this interpretation comes from its graded sensitivity to orthographic overlap between prime and target in masked priming; N250 amplitude is most attenuated for full identity repetitions, intermediate for primes sharing all but one letter, and largest for fully unrelated pairs \citep{holcombTimeCourseVisual2006, holcomb_exploring_2007}. The N250 is also sensitive to sublexical phonological overlap, with \citet{graingerTimeCourseOrthographic2006} identifying separable early orthographic and later phonological subcomponents, consistent with dual-route architectures of word recognition.
	
	The N400 is a negative-going component peaking around 400 ms, with a broad centro-parietal distribution \citep{kutasReadingSenselessSentences1980}. Across more than four decades of research, its amplitude has been shown to vary with semantic priming, cloze probability, lexical association, concreteness, and semantic richness, while remaining insensitive to physical and syntactic manipulations that leave meaning intact \citep{kutasThirtyYearsCounting2011}. Critically for the present study, the N400 is reliably modulated in single-word and word-pair paradigms, with larger amplitudes for semantically unrelated than related targets, making it an appropriate real-time measure of lexical-semantic access.
	
	Although both components are negative-going deflections occurring in adjacent time windows, they are dissociable by topography—the N250 has a more frontal distribution, the N400 a posterior one \citep{holcombTimeCourseVisual2006}—and, more decisively, by a double dissociation across priming types. Lexical status selectively modulates the N400 in that priming effects are restricted to real words with semantic representations and do not extend to pseudowords \citep{kiyonagaMaskedCrossmodalRepetition2007}, whereas the N250 is equally modulated by words and pseudowords alike. The reverse pattern holds for semantic relatedness; while associative priming reliably modulates the N400, it produces no N250 effect \citep{graingerWatchingWordGo2009}. This double dissociation—pseudoword priming engaging the N250 but not the N400, semantic priming engaging the N400 but not the N250—confirms that the two components index qualitatively distinct processing stages, and rules out the interpretation of the N250 as merely an early onset of the N400.
	
	
	\subsection{Aims and Predictions}
	
	The present study has two complementary aims. The first is to establish whether base frequency and morphological family size exert their effects at distinct temporal stages of visual word recognition, using ERP to track processing in real time. The second, and primary, aim builds on this: to examine whether individual differences in orthographic and semantic processing, captured by two continuous dimensions, Orthographic Sensitivity (OS) and Language Proficiency (LP), derived from a principal components analysis of spelling ability, vocabulary knowledge, and print exposure, modulate these effects selectively, at the processing stages where each variable is predicted to operate.
	
	Although OS and LP are orthogonal dimensions rather than poles of a single continuum, they nonetheless differ in their primary loading: OS captures relatively greater precision of sublexical orthographic representation, while LP captures relatively greater breadth of lexical-semantic knowledge. The core prediction is therefore a double dissociation—Orthographic Sensitivity should modulate base frequency effects at the level of early form-based processing, as indexed by the N250, while Language Proficiency should modulate family size effects at the level of semantic access, as indexed by the N400. This prediction follows directly from \citeauthor{andrewsMorphologicalPrimingStronger2013}'s \citeyearpar{andrewsMorphologicalPrimingStronger2013} finding that the balance between orthographic and semantic processing is a genuine and measurable axis of individual variation, combined with the broadly shared theoretical assumption that base frequency reflects form-level access and family size reflects the richness of semantic representation. If this double dissociation is observed, it would constitute direct electrophysiological evidence that the distinction between form-based and meaning-based morphological processing is not merely a property of competing computational models but a real axis of variation in the reading system. If, instead, both variables are modulated by the same individual difference dimension at the same processing stage, this would suggest that the two effects are less dissociable than current multi-stage accounts propose.
	
	The first analysis examines these predictions in real word recognition. Because real words were not matched across morphological variables, reflecting the naturalistic distributional properties of morphologically complex words rather than an orthogonally controlled factorial design, base frequency, family size, and morphological complexity are treated as continuous predictors in a regression framework for reaction times. For ERP data, items were median-split into high and low groups separately for base frequency and family size to yield stable waveforms for comparison. The two measures therefore carry somewhat different inferential weight—the RT regression speaks to the independent continuous contributions of each variable, while the ERP analyses speak to the direction and timing of effects at the extremes of each variable's distribution.
	
	The second analysis examines whether these individual difference modulations generalise to novel morphological structure, using morphologically complex nonwords. Morphologically structured nonwords (e.g., \textit{gasful}: real stem + real affix) were compared against matched controls (e.g., \textit{gasfil}: real stem + pseudoaffix), holding stem identity constant while manipulating suffix legitimacy. Because these items have no stored whole-word representations, they place maximal pressure on the system's generalisation mechanism and allow stem-level variables to be examined independently of whole-word familiarity. The key prediction is that Orthographic Sensitivity should modulate early N250 responses to the suffix contrast. For nonwords specifically, no stored whole-word entry exists, so early processing must rely on readers' sublexical orthographic knowledge—readers with more precise representations of real suffix forms should show greater N250 sensitivity to the distinction between real and pseudo-suffixes. Language Proficiency, by contrast, should predict later N400 modulation: readers with broader and more richly specified lexical-semantic knowledge should show stronger N400 sensitivity to whether a nonword's stem belongs to a large or productive morphological family, since it is this family membership that determines the richness of the semantic activation the form can support. However, the predicted dissociation between OS at the N250 and LP at the N400 may not be fully clean for nonwords. Because readers' familiarity with real suffix forms is acquired partly through exposure to morphological families—readers learn that \emph{-ness} is a real suffix in part because they have encountered many \emph{-ness} words with related meanings—the breadth of a reader's morphological vocabulary, which contributes to Language Proficiency, may also enrich early processing of novel suffix forms. This leads to the weaker prediction that LP may contribute to N250 modulation for nonwords alongside OS.
	
	\section{Methods}
	
	\subsection{Materials}
	
	The experimental stimulus set included 200 nonwords, constructed to manipulate morphological structure while controlling surface form. In the \textit{complex nonword} condition, we combined existing English base words (e.g., \textit{gas}) with existing derivational suffixes (e.g., \textit{gasful}), to yield 100 morphologically plausible complex nonwords. Each combination was syntactically legal and reflected common English derivational patterns. A total of 16 suffixes were used, each attached to four different stems.
	
	In the \textit{simple nonword} condition, we combined the same base words as in the complex non-word condition with non-suffix endings that were orthographically similar to the real suffixes used in the complex condition but did not constitute valid morphemes, to yield 100 orthographically plausible simple non-words. These pseudoaffix endings were created by changing the central letter of each suffix (e.g., \textit{gasful}–\textit{gasfil}), preserving length and orthographic neighborhood characteristics while eliminating morphological structure. Thus each complex non-word had an orthographically matched simple counterpart.
	
	Because the same base appeared in both the simple and complex conditions, the experimental nonwords were distributed across two lists such that no participant encountered the same base more than once. In addition to the experimental nonwords, the stimulus set included 100 monomorphemic fillers (50 real words and 50 nonwords) to balance lexical status and reduce predictability. The experimental stimuli were drawn from \citet{sawiIndividualDifferencesSensitivity2016}
	
	Estimates of cumulative base (root) frequency, surface frequency, and morphological family size were obtained from the CELEX lexical database. Morphological family size was calculated by identifying all morphologically related forms sharing the same base (e.g., \textit{calculate}, \textit{calculable}, \textit{calculation}, \textit{calculator}). Following prior work \citep[e.g.,][]{dejongMorphologicalFamilySize2000, dejongProcessingRepresentationDutch2002}, both compounds (e.g., \textit{WATCHTOWER}) and hyphenated compounds (e.g., \textit{CHECK-IN}) were included, as they contribute to morphological connectivity within the lexicon. Base frequency was computed by summing the lemma frequencies of all confirmed family members.
	
	For reaction time analyses, all lexical and individual-difference variables were treated as continuous predictors to preserve information and maximize statistical power. In contrast, ERP analyses were conducted on condition-averaged waveforms to improve signal-to-noise ratio; as a result, item-level predictors such as base frequency and family size were operationalized as categorical condition factors rather than continuous trial-level covariates. This distinction reflects standard practice in ERP research and aligns the statistical treatment of predictors with the structure of the underlying data.
	
	\subsection{Individual Differences and Grouping Variables}
	
	To quantify individual differences in linguistic experience and orthographic knowledge, participants completed three standardized measures—the Nelson–Denny Vocabulary Test \citep{brownNelsonDennyReadingTest1993}, a spelling task adapted from \citet{sawiIndividualDifferencesSensitivity2016}, and the Author Recognition Test (ART; \citealt{stanovichExposurePrintOrthographic1989}). Scores on these measures were moderately correlated ($r$s = .29–.51, all $p < .02$), indicating partially overlapping but dissociable components of lexical experience.
	
	To capture the shared and distinct variance among these measures, we conducted a principal components analysis (PCA) on standardized (z-scored) vocabulary, spelling, and ART scores. PCA was selected to derive orthogonal dimensions of individual differences without imposing a priori weighting or categorical group boundaries. Two components had eigenvalues greater than 1, jointly accounting for 83.5\% of the total variance (58.0\% for Dimension~1 and 25.5\% for Dimension~2).
	
	Component loadings (Table~\ref{tab:pca_loadings}) indicated that Dimension~1 loaded most strongly on vocabulary and print exposure (ART), with a secondary contribution from spelling accuracy, and was interpreted as indexing \textit{Language Proficiency}. Dimension~2 loaded most strongly on spelling accuracy, with weaker and oppositely signed contributions from vocabulary and ART, and was interpreted as indexing \textit{Orthographic Sensitivity}:
	
	\begin{itemize}
		\item[-] \textit{Dimension~1 (Language Proficiency)}: vocabulary (.82), ART (.81), spelling (.65)
		\item[-] \textit{Dimension~2 (Orthographic Sensitivity)}: spelling (.76), vocabulary ($- .29$), ART ($- .32$)
	\end{itemize}
	
	The two components were statistically independent ($r \approx 0$), supporting their interpretation as orthogonal dimensions of reading-related skill. These continuous component scores were used as predictors in all behavioral and ERP analyses.
	
	
	\begin{table}[ht]
		\centering
		\caption{PCA loadings and variance explained for the three lexical-experience measures.}
		\label{tab:pca_loadings}
		\begin{tabular}{lcc}
			\toprule
			\textbf{Measure} & \textbf{\shortstack{Dimension~1 \\ (Language Proficiency)}} & \textbf{\shortstack{Dimension~2 \\ (Orthographic Sensitivity)}} \\
			\midrule
			Vocabulary (z\_vcb) & 0.82 & $-0.29$ \\
			Spelling (z\_spl) & 0.65 & 0.76 \\
			Author Recognition Test (z\_art) & 0.81 & $-0.32$ \\
			\midrule
			\textit{Eigenvalue} & 1.74 & 0.76 \\
			\textit{Variance explained (\%)} & 58.0 & 25.5 \\
			\bottomrule
		\end{tabular}
	\end{table}
	
	\section{Statistical Analyses}
	
	The analyses are organized around the study’s primary theoretical aim—determining whether Orthographic Sensitivity and Language Proficiency dissociably modulate the timing and nature of morphological processing. We predicted that Orthographic Sensitivity would selectively modulate early, form-based processing indexed by the N250, while Language Proficiency would selectively modulate later semantic integration indexed by the N400. The behavioral and ERP results are reported separately for words and nonwords in each section, followed by a summary that evaluates these predictions against the obtained pattern.
	
	\subsection{Behavioral Data (Reaction Times)}
	Reaction time (RT) analyses were conducted on correct responses only. RTs were log-transformed to reduce positive skew and to better meet model assumptions of normality and homoscedasticity. Separate analyses were conducted for word and nonword trials: for words, base frequency and family size were the primary morphological predictors; for nonwords, morphological complexity (real vs. pseudo-suffix) was the key manipulation, with base frequency and family size included as covariates (see Aims and Predictions).
	
	RTs were analyzed using linear mixed-effects models fit with the \texttt{lme4} package \citep{batesFittingLinearMixedeffects2015} in R \citep{rcoreteamLanguageEnvironmentStatistical2026}. All models included random intercepts for participants and items, and random slopes for within-participant predictors where supported by the data. Continuous lexical predictors (log base frequency and log family size) were mean-centered and scaled prior to analysis. Individual-difference measures—Orthographic Sensitivity and Language Proficiency—were entered as continuous predictors corresponding to standardized PCA scores.
	
	For word trials, models tested the effects of base frequency, family size, Orthographic Sensitivity, and Language Proficiency, as well as interactions between Orthographic Sensitivity and lexical familiarity measures. For nonword trials, models additionally included morphological complexity (simple vs. complex) as a categorical predictor and tested whether complexity effects were modulated by individual differences. Treatment (reference) coding was used for categorical predictors in RT analyses. Model significance was assessed using Type III tests with Satterthwaite-approximated degrees of freedom. Fixed-effect estimates are reported alongside model-based estimated marginal means where relevant. Means are reported in milliseconds following back-transformation from the log scale used in analysis.
	
	\subsection{ERP Data}
	For word stimuli, ERP trials were averaged separately according to median splits on base frequency and morphological family size, enabling separate analyses of each variable at the N250 and N400. For nonword stimuli, trials were averaged according to a median split on morphological family size. Family size served different roles across the two time windows: at the N400 it was a primary predictor of theoretical interest, entering into the predicted Complexity $\times$ Family Size $\times$ LP interaction; at the N250 it was retained as a covariate to reduce item-level variance. Because base frequency and family size are properties of the stem shared across both members of each complexity pair (e.g., \textit{gasful} and \textit{gasfil} share the same stem), neither variable could confound the complexity manipulation. Family size was preferred over base frequency as the N250 covariate on empirical grounds: preliminary model comparisons indicated that family size accounted for substantially more item-level variance in the nonword N250 data (residual variance: 4.07 vs.\ 27.47). The pattern of interactions reported below was equivalent in direction when base frequency was used as the covariate instead.
	
	ERP analyses were conducted on mean amplitudes computed from condition-averaged waveforms. Separate datasets were constructed for words and nonwords, and analyses were conducted independently for the N250 and N400 time windows. Because ERP amplitudes were derived from condition-level averages rather than trial-level data, morphological predictors in ERP analyses were treated as categorical condition factors rather than continuous covariates.
	
	ERP amplitudes were analysed using linear mixed-effects models. For word stimuli, separate models were fit for base frequency and family size, each entered as the primary morphological predictor along with Orthographic Sensitivity, Language Proficiency, and their interaction with the morphological predictor. For nonword stimuli, model structure differed across time windows to reflect the distinct role of family size at each stage. At the N250, models included fixed effects of Morphological Complexity, Orthographic Sensitivity, Language Proficiency, interactions between Complexity and each individual difference dimension, and Morphological Family Size as a covariate. At the N400, models included fixed effects of Morphological Complexity, Morphological Family Size, Language Proficiency, their three-way interaction, and Orthographic Sensitivity as a covariate, reflecting the theoretical prediction that LP selectively modulates semantic integration of morphological family information at this later stage. All models included random intercepts for participants and for electrode site nested within participant, allowing baseline amplitude differences across electrodes to be modelled while preserving the spatial structure of the data.
	
	Categorical predictors in ERP analyses were effect-coded (sum-to-zero coding), such that model intercepts represent grand mean amplitudes and main effects reflect average effects across conditions. Orthographic Sensitivity and Language Proficiency were entered as continuous, mean-centred predictors. Time window was not included as a factor in the models; instead, separate models were fit for the N250 and N400 to preserve component-specific interpretation.
	
	Statistical significance of fixed effects was assessed using Type III $F$-tests with Satterthwaite-approximated degrees of freedom as implemented in the \texttt{lmerTest} package [CITATION]. Where higher-order interactions involving individual-difference measures were significant, follow-up analyses were conducted using estimated marginal means (\texttt{emmeans}). Interactions involving Orthographic Sensitivity and Language Proficiency were probed by estimating contrasts at $\pm$1 SD of each dimension, allowing direct tests of whether morphological effects differed as a function of reader skill.
	
	\subsection{Model Estimation and Inference}
	All mixed-effects models were fit using restricted maximum likelihood estimation. To ensure stable convergence, models were estimated using the BOBYQA optimizer. Continuous predictors were centered to facilitate interpretation of interaction terms and model intercepts. For predictors with one degree of freedom, F-tests from the analysis-of-variance tables are mathematically equivalent to squared t-tests from model summaries; all inferential conclusions are therefore invariant to the choice of reporting format. Across all analyses, inferential emphasis was placed on theoretically motivated interactions rather than exhaustive post hoc comparisons. Follow-up contrasts were conducted only where warranted by significant interactions and were limited to tests directly addressing the study’s hypotheses. Estimated marginal means and confidence intervals were obtained using the \texttt{emmeans} package [CITATION]; confidence intervals are based on asymptotic approximation, which is appropriate given the sample size.
	
	\section{Results}
	
	The analyses are organized around two complementary sources of evidence bearing on the study's primary theoretical aim—determining whether Orthographic Sensitivity and Language Proficiency dissociably modulate the timing and nature of morphological processing. Reaction time models included individual difference measures as main effects only, reflecting the theoretical framing of behavioral responses as indices of overall processing efficiency. Interactions between morphological variables and individual differences were examined in the ERP models, where the temporal resolution of the N250 and N400 allows these effects to be attributed to specific processing stages. Results are reported separately for words and nonwords. For each stimulus type, behavioral results are presented first, followed by ERP results organized by component, N250 then N400, reflecting the temporal sequence of morphological processing. This organization reflects the distinct theoretical roles of the two stimulus types—the word analyses examine the time course of base frequency and family size effects in stored lexical representations, while the nonword analyses examine whether morphological sensitivity generalizes to novel forms and whether individual differences modulate that generalization. Within each section, the pattern of effects is evaluated against the predictions derived from structured parsing accounts and the discriminative learning framework, with particular attention to whether Orthographic Sensitivity selectively modulates early form-based processing and Language Proficiency selectively modulates later semantic integration. The specific fixed effects, interactions, and random effects structure for each model are summarised in Table~\ref{tab:models}, followed by a brief explanation of the asymmetric interaction structure.
	
	Table~\ref{tab:models} summarises the fixed effects and interaction terms included in each model. ERP models for words and nonwords tested the interaction between morphological variables and the individual difference dimension predicted to be most relevant at each processing stage, Orthographic Sensitivity at the N250, reflecting early form-based processing, and Language Proficiency at the N400, reflecting later semantic integration. Reaction time models included individual difference measures as main effects only, reflecting the theoretical framing of behavioral responses as indices of overall processing efficiency rather than processing architecture.
	
	\begin{table}[h!]
		\centering
		\small
		\renewcommand{\arraystretch}{1.4}
		\caption{Summary of linear mixed-effects models. OS = Orthographic Sensitivity; LP = Language Proficiency; BF = Base Frequency; FS = Family Size. All models fit by REML using the \texttt{bobyqa} optimizer. Degrees of freedom estimated using Satterthwaite's method. RE = random effects; Subj = participant; Elec = electrode site.}
		\label{tab:models}
		\begin{tabular}{p{1.6cm} p{1.2cm} p{1.3cm} p{3.5cm} p{2.8cm} p{2.5cm}}
			\toprule
			\textbf{Stimulus} & \textbf{Outcome} & \textbf{Variable} & \textbf{Fixed effects} & \textbf{Interaction(s)} & \textbf{RE structure} \\
			\midrule
			Words  & RT  & BF + FS & BF, FS, OS, LP       & None                & Subj, Item \\
			Words  & N250 & BF   & BF, OS, LP         & BF $\times$ OS           & Subj, Elec/Subj \\
			Words  & N250 & FS   & FS, OS, LP         & FS $\times$ OS           & Subj, Elec/Subj \\
			Words  & N400 & BF   & BF, LP, OS         & BF $\times$ LP           & Subj, Elec/Subj \\
			Words  & N400 & FS   & FS, LP, OS         & FS $\times$ LP           & Subj, Elec/Subj \\
			Nonwords & RT  & BF + FS & Complexity, BF, FS, OS, LP & None                & Subj (Complexity slope), Item \\
			Nonwords & N250 & FS   & Complexity, OS, LP, FS   & Cmplx $\times$ OS, Cmplx $\times$ LP & Subj, Elec/Subj \\
			Nonwords & N400 & FS   & Complexity, FS, LP, OS   & Cmplx $\times$ FS $\times$ LP    & Subj, Elec/Subj \\
			\bottomrule
		\end{tabular}
	\end{table}
	
	The asymmetric interaction structure reflected in Table~\ref{tab:models} warrants brief explanation. For word stimuli, separate ERP models were run for base frequency and family size at each time window, reflecting the fact that these variables were operationalised as separate median-split factors in distinct datasets. At the N250, both the base frequency and family size models tested the interaction with Orthographic Sensitivity, consistent with the prediction that early form-based processing is selectively modulated by sublexical orthographic knowledge; Language Proficiency was included as a covariate in both models. The corresponding N400 models reversed this structure, testing interactions with Language Proficiency and including Orthographic Sensitivity as a covariate.
	
	For nonword stimuli, base frequency was included as a covariate in the RT model but was not independently examined in the ERP analyses. At the N250, morphological family size was included as a covariate to reduce item-level variance rather than as a predictor of theoretical interest; the primary fixed effects were Morphological Complexity and its interactions with both individual difference dimensions, reflecting the prediction that early processing of novel forms is modulated by reader skill. At the N400, family size was elevated to a primary predictor, entering into the three-way Complexity $\times$ Family Size $\times$ LP interaction that constitutes the central test of this time window; Orthographic Sensitivity was included as a covariate. This distinction in the role of family size across the two nonword ERP models reflects the theoretical rationale developed in the Aims and Predictions section: family size indexes the richness of the morphological neighbourhood and is predicted to modulate semantic integration specifically at the N400, where LP selectively drives processing. This structure ensures that each model is specified to answer a theoretically motivated question rather than exhaustively testing all possible interactions.
	
	\subsection{Words}
	
	\paragraph*{Words RT}
	Reaction times (RTs) for correct lexical decisions to word stimuli were analyzed using a linear mixed-effects model with fixed effects of Base Frequency, Family Size, Orthographic Sensitivity (OS), and Language Proficiency (LP). The model included random intercepts for participants and items. The analysis revealed robust main effects of item-level lexical familiarity. As base frequency increased, recognition times decreased, $\beta = -0.017$, $SE = 0.005$, $t(100.97) = -3.66$, $p < .001$, $\eta^2_p = .12$. Similarly, words from larger morphological families elicited faster responses than those from smaller families, $\beta = -0.013$, $SE = 0.005$, $t(93.59) = -2.64$, $p = .010$, $\eta^2_p = .07$. Orthographic Sensitivity showed a marginal trend toward faster overall response times, $\beta = -0.035$, $SE = 0.019$, $t(63.84) = -1.85$, $p = .069$, $\eta^2_p = .05$. Language Proficiency did not exert a reliable main effect on RTs, $\beta = 0.012$, $SE = 0.012$, $t(63.97) = 0.94$, $p = .35$. The model explained 3.3\% of the variance via fixed effects ($R^2_{\mathrm{marginal}} = .033$) and 41.7\% overall ($R^2_{\mathrm{conditional}} = .417$). These findings indicate that word recognition speed is driven primarily by item-level lexical familiarity; Orthographic Sensitivity shows only a marginal and unstable influence on overall response efficiency and does not selectively modulate morphological effects on response times.
	
	\paragraph*{Words N250: Base Frequency}
	Mean N250 amplitudes for word stimuli were analyzed using a linear mixed-effects model with fixed effects of Base Frequency (high vs.\ low), Orthographic Sensitivity (OS), Language Proficiency (LP), and the interaction between Base Frequency and OS. The model included random intercepts for participants and for electrode site nested within participant. The analysis revealed a significant main effect of Base Frequency—words from higher-frequency bases elicited more negative N250 amplitudes than words from lower-frequency bases, $M_{\mathrm{high}} = -1.63\,\mu$V, 95\% CI $[-2.56, -0.69]$, $M_{\mathrm{low}} = -1.29\,\mu$V, 95\% CI $[-2.22, -0.36]$, $\beta = -0.318$, $SE = 0.100$, $t(1618) = -3.18$, $p = .002$, $\eta^2_p = .006$. Neither Orthographic Sensitivity, $\beta = 0.534$, $SE = 0.520$, $t(57) = 1.03$, $p = .309$, nor Language Proficiency, $\beta = -0.536$, $SE = 0.378$, $t(57) = -1.42$, $p = .162$, yielded reliable main effects. Critically, the interaction between Base Frequency and Orthographic Sensitivity was significant, $\beta = -0.455$, $SE = 0.112$, $t(1618) = -4.06$, $p < .001$, $\eta^2_p = .010$, indicating that the magnitude of early base frequency effects is modulated by orthographic skill. The negative direction of the interaction indicates that readers higher in Orthographic Sensitivity showed an amplified N250 base frequency effect—a larger difference in early processing amplitude between high- and low-frequency stems—relative to readers lower in OS, consistent with the interpretation that more precise sublexical orthographic representations strengthen early morpheme-level activation signatures. The model explained 1.8\% of variance via fixed effects ($R^2_{\mathrm{marginal}} = .018$) and 80.1\% overall ($R^2_{\mathrm{conditional}} = .801$).
	
	\paragraph*{Words N250: Family Size}
	Mean N250 amplitudes for word stimuli were analyzed using a linear mixed-effects model with fixed effects of Family Size, Orthographic Sensitivity (OS), Language Proficiency (LP), and the interaction between Family Size and OS. The model included random intercepts for participants and for electrode site nested within participant. The analysis revealed a robust main effect of Family Size—words from smaller morphological families elicited more negative N250 amplitudes than words from larger families, $M_{\mathrm{large}} = -0.57\,\mu$V, 95\% CI $[-1.03, -0.11]$, $M_{\mathrm{small}} = -0.90\,\mu$V, 95\% CI $[-1.36, -0.44]$, $\beta = 0.324$, $SE = 0.058$, $t(1618) = 5.54$, $p < .001$, $\eta^2_p = .02$. Neither Orthographic Sensitivity, $\beta = 0.261$, $SE = 0.256$, $t(57) = 1.02$, $p = .31$, nor Language Proficiency, $\beta = -0.279$, $SE = 0.187$, $t(57) = -1.50$, $p = .14$, yielded reliable main effects. The interaction between Family Size and Orthographic Sensitivity did not reach significance but showed a marginal trend, $\beta = -0.111$, $SE = 0.065$, $t(1618) = -1.70$, $p = .089$, $\eta^2_p = .002$, suggesting that the family size effect at the N250 may be modulated by orthographic skill, though this requires cautious interpretation given the marginal status of the effect. The model explained 1.9\% of the variance in N250 amplitude via fixed effects ($R^2_{\mathrm{marginal}} = .019$) and 73.8\% overall ($R^2_{\mathrm{conditional}} = .738$), with the large conditional $R^2$ reflecting substantial variability in N250 amplitude across participants and electrode sites. Together, these results indicate that morphological family structure robustly modulates early morpho-orthographic processing for words, with limited but suggestive evidence that orthographic skill may shape the sensitivity of this effect.
	
	\paragraph*{Words N400: Base Frequency}
	
	Mean N400 amplitudes for word stimuli were analysed using a linear  mixed-effects model with fixed effects of Base Frequency (high vs.\  low), Language Proficiency (LP), Orthographic Sensitivity (OS), and  the interaction between Base Frequency and LP. The model included  random intercepts for participants and for electrode site nested within  participant. The analysis revealed a significant main effect of Base  Frequency --- averaged across LP levels, words from higher-frequency  bases elicited more positive N400 amplitudes than words from  lower-frequency bases, $M_{\mathrm{high}} = 0.71\,\mu$V, 95\% CI  $[-0.44, 1.86]$, $M_{\mathrm{low}} = 0.52\,\mu$V, 95\% CI $[-0.63,  1.66]$, $\beta = 0.214$, $SE = 0.106$, $t(1618) = 2.01$, $p = .044$,  $\eta^2_p = .003$. A significant main effect of Orthographic  Sensitivity was also observed, $\beta = 1.290$, $SE = 0.640$, $t(57)  = 2.02$, $p = .049$, $\eta^2_p = .07$, reflecting overall differences  in N400 amplitude as a function of Orthographic Sensitivity,  independently of morphological structure --- a pattern consistent with  the OS main effect observed in the family size model below. Although  this effect was significant, the $p$-value was close to the  conventional threshold ($p = .049$), and should be interpreted with  appropriate caution. Language Proficiency did not yield a reliable main  effect, $\beta = 0.047$, $SE = 0.465$, $t(57) = 0.10$, $p = .920$.
	
	Critically, the interaction between Base Frequency and Language  Proficiency was significant, $\beta = 0.344$, $SE = 0.087$, $t(1618)  = 3.98$, $p < .001$, $\eta^2_p = .010$. Follow-up contrasts examined  the base frequency effect across the observed LP range ($-2.7$ to  $+3.1$ SD). The direction of the base frequency effect reversed across  the LP continuum. At LP $= -2$ SD, high base frequency words elicited  more negative N400 amplitudes than low base frequency words ($\Delta =  -0.470\,\mu$V, $t(1618) = -2.38$, $p = .018$; EMM: High $=  0.28\,\mu$V, Low $= 0.76\,\mu$V). At LP $= -1$ SD the contrast was  not significant ($\Delta = -0.130\,\mu$V, $t(1618) = -0.97$, $p =  .333$). At mean LP the direction had reversed, with high base frequency  words now eliciting more positive N400 amplitudes than low base  frequency words ($\Delta = 0.210\,\mu$V, $t(1618) = 2.01$, $p =  .044$; EMM: High $= 0.72\,\mu$V, Low $= 0.51\,\mu$V), and this  pattern strengthened monotonically at LP $= +1$ SD ($\Delta =  0.560\,\mu$V, $t(1618) = 3.99$, $p < .001$; EMM: High $= 0.94\,\mu$V,  Low $= 0.38\,\mu$V) and LP $= +2$ SD ($\Delta = 0.900\,\mu$V,  $t(1618) = 4.37$, $p < .001$; EMM: High $= 1.16\,\mu$V, Low $=  0.26\,\mu$V). The participant LP scores ranged from $-2.67$ to $+3.11$,  confirming that estimates at $\pm 2$ SD are grounded in observed data  rather than extrapolation. Together, these contrasts indicate that LP  modulates the contribution of base frequency to N400 amplitude in a  manner that crosses over across the LP continuum, with the nature of  that modulation shifting from a pattern in which high base frequency  words are harder to process semantically at low LP, to one in which  they are easier to process at high LP. The model explained 2.6\% of  variance via fixed effects ($R^2_{\mathrm{marginal}} = .026$) and  82.5\% overall ($R^2_{\mathrm{conditional}} = .825$).
	
	\paragraph*{Words N400: Family Size}
	Mean N400 amplitudes for word stimuli were analyzed using a linear mixed-effects model with fixed effects of Family Size, Language Proficiency (LP), Orthographic Sensitivity (OS), and the interaction between Family Size and LP. The model included random intercepts for participants and for electrode site nested within participant. The analysis revealed a robust main effect of Family Size—words from smaller morphological families elicited greater N400 amplitudes than words from larger families, $M_{\mathrm{large}} = 0.518\,\mu$V, 95\% CI $[-0.054, 1.089]$, $M_{\mathrm{small}} = 0.062\,\mu$V, 95\% CI $[-0.509, 0.634]$, $\beta = 0.452$, $SE = 0.064$, $t(1618) = 7.07$, $p < .001$, , $\eta^2_p = .03$, indicating facilitated semantic integration for words supported by richer morphological families. A significant main effect of Orthographic Sensitivity was also observed, $\beta = 0.677$, $SE = 0.318$, $t(57) = 2.13$, $p = .038$, $\eta^2_p = .07$, reflecting overall differences in N400 amplitude as a function of orthographic skill independently of morphological structure. Language Proficiency did not yield a reliable main effect, $\beta = 0.001$, $SE = 0.231$, $t(57) = 0.005$, $p = .996$, nor did the Family Size $\times$ LP interaction reach significance, $\beta = -0.071$, $SE = 0.052$, $t(1618) = -1.37$, $p = .172$. These results indicate that semantic integration of morphological family information for familiar words is robust and not selectively modulated by individual differences in language proficiency.
	
	
	\subsection{Nonwords}
	
	\paragraph*{Nonwords RT}
	Reaction times (RTs) for correct lexical decisions to nonwords were analyzed using a linear mixed-effects model with fixed effects of Morphological Complexity (simple vs.\ complex), Base Frequency, Family Size, Orthographic Sensitivity (OS), and Language Proficiency (LP). The model included random intercepts for participants and items and by-participant random slopes for Complexity. The analysis revealed a strong main effect of Morphological Complexity—complex nonwords elicited slower responses than simple nonwords, $M_{\mathrm{complex}} = 718.9$ ms, 95\% CI $[695.2, 743.4]$, $M_{\mathrm{simple}} = 684.6$ ms, 95\% CI $[661.5, 708.5]$, $\beta = 0.049$, $SE = 0.005$, $t(62.68) = 9.13$, $p < .001$, $\eta^2_p = .57$. In addition, Base Frequency exerted a reliable inhibitory effect, $\beta = 0.015$, $SE = 0.005$, $t(98.21) = 3.05$, $p = .003$, $\eta^2_p = .09$, such that nonwords derived from higher-frequency bases were responded to more slowly than those from lower-frequency bases. Family Size did not significantly influence response times, $\beta = -0.006$, $SE = 0.005$, $t(96.64) = -1.24$, $p = .22$. Orthographic Sensitivity showed a significant main effect, with higher-sensitivity participants responding more quickly overall, $\beta = -0.041$, $SE = 0.018$, $t(62.21) = -2.23$, $p = .030$, $\eta^2_p = .07$. In contrast, there was no effect of Language Proficiency on nonword RTs, $\beta = 0.003$, $SE = 0.012$, $t(62.91) = 0.28$, $p = .78$. The model explained 5.0\% of the variance via fixed effects ($R^2_{\mathrm{marginal}} = .050$) and 49.4\% overall ($R^2_{\mathrm{conditional}} = .494$). Taken together, these results indicate that nonword decision latencies are strongly shaped by morphological complexity and item-level lexical familiarity, particularly base frequency, while individual differences primarily affect overall response efficiency, with more orthographically sensitive readers responding more quickly across conditions.
	
	
	\paragraph*{Nonwords N250}
	
	Mean N250 amplitudes for nonword stimuli were analysed using a linear mixed-effects model with fixed effects of Morphological Complexity, Orthographic Sensitivity (OS), Language Proficiency (LP), interactions between Complexity and each individual difference dimension, and Morphological Family Size as a covariate. The model included random intercepts for participants and for electrode site nested within participant, and accounted for 57.5\% of variance in N250 amplitude when random effects were included ($R^2_{\mathrm{conditional}} = .575$, $R^2_{\mathrm{marginal}} = .025$).
	
	The analysis revealed a significant main effect of Morphological Complexity, with complex nonwords eliciting more negative N250 amplitudes than simple nonwords overall ($M_{\text{complex}} = -0.759\,\mu$V $[-1.148, -0.371]$, $M_{\text{simple}} = -0.526\,\mu$V $[-0.914, -0.138]$), $\beta = -0.233$, $SE = 0.050$, $t(4856) = -4.63$, $p < .001$. Morphological Family Size, included as a covariate, yielded a significant effect, $\beta = 0.182$, $SE = 0.050$, $t(4856) = 3.63$, $p < .001$, with large-family nonwords eliciting less negative amplitudes than small-family nonwords. Neither OS nor LP yielded reliable main effects, $\beta = 0.298$, $SE = 0.220$, $t(57) = 1.35$, $p = .181$; $\beta = -0.207$, $SE = 0.160$, $t(57) = -1.29$, $p = .202$, respectively.
	
	Critically, Morphological Complexity interacted significantly with both individual difference dimensions. The Complexity $\times$ OS interaction was significant, $\beta = 0.392$, $SE = 0.056$, $t(4856) = 6.98$, $p < .001$, $\eta^2_p = .010$. Follow-up contrasts examined the complexity effect across the observed OS range ($-1.9$ to $+1.9$ SD). The complexity effect was large and reliable at OS $= -1.9$ SD ($\Delta = 0.990\,\mu$V, $t = 8.29$, $p < .001$) and at OS $= -1$ SD ($\Delta = 0.642\,\mu$V, $t = 8.33$, $p < .001$), attenuated but still significant at mean OS ($\Delta = 0.250\,\mu$V, $t = 4.98$, $p < .001$), and reversed in direction but only marginally significant at OS $= +1$ SD ($\Delta = -0.142\,\mu$V, $t = -1.93$, $p = .053$). At the upper extreme of the observed OS range (OS $= +1.9$ SD) the reversal was significant ($\Delta = -0.490\,\mu$V, $t = -4.27$, $p < .001$). Examination of estimated marginal means indicated that this reversal was driven primarily by the trajectory of complex nonwords, which decreased in negativity across the OS range (EMM: $-1.720\,\mu$V at OS $= -1.9$ SD to $+0.160\,\mu$V at OS $= +1.9$ SD), while simple nonwords also decreased in negativity but more slowly (EMM: $-0.720\,\mu$V at OS $= -1.9$ SD to $-0.340\,\mu$V at OS $= +1.9$ SD).
	
	The Complexity $\times$ LP interaction was also significant, $\beta = 0.358$, $SE = 0.041$, $t(4856) = 8.76$, $p < .001$, $\eta^2_p = .016$. Follow-up contrasts examined the complexity effect across the observed LP range ($-2.7$ to $+3.1$ SD). The complexity effect was large and reliable at LP $= -2$ SD ($\Delta = 0.930\,\mu$V, $t = 9.89$, $p < .001$) and at LP $= -1$ SD ($\Delta = 0.574\,\mu$V, $t = 9.05$, $p < .001$), attenuated but still significant at mean LP ($\Delta = 0.216\,\mu$V, $t = 4.31$, $p < .001$), and significantly reversed at LP $= +1$ SD ($\Delta = -0.141\,\mu$V, $t = -2.14$, $p = .032$) and LP $= +2$ SD ($\Delta = -0.500\,\mu$V, $t = -5.12$, $p < .001$). In contrast to the OS pattern, examination of estimated marginal means indicated that the LP reversal was driven entirely by the trajectory of simple nonwords, which became increasingly negative across the LP range (EMM: $+0.230\,\mu$V at LP $= -2$ SD to $-1.320\,\mu$V at LP $= +2$ SD), while complex nonwords remained essentially flat (EMM: $-0.700\,\mu$V at LP $= -2$ SD to $-0.820\,\mu$V at LP $= +2$ SD). Model fit indices were $R^2_{\text{marginal}} = .025$ and $R^2_{\text{conditional}} = .575$.
	
	
	
	\paragraph*{Nonwords N400}
	Mean N400 amplitudes for nonword stimuli were analysed using a linear mixed-effects model with fixed effects of Morphological Complexity, Family Size, Language Proficiency (LP), Orthographic Sensitivity (OS), and all interactions among Complexity, Family Size, and LP. The model included random intercepts for participants and for electrode site nested within participant.
	
	The analysis revealed a significant main effect of Morphological Complexity, a significant Complexity $\times$ Family Size interaction, a Complexity $\times$ LP interaction, a Family Size $\times$ LP interaction, and critically, a significant three-way Complexity $\times$ Family Size $\times$ LP interaction. Main effects and lower-order interactions are reported as context for the three-way, which is the primary theoretical finding.
	
	Morphological Complexity yielded a robust main effect---complex nonwords elicited more negative N400 amplitudes than simple nonwords, $\beta = -0.460$, $SE = 0.061$, $t(4854) = -7.48$, $p < .001$, $\eta^2_p = .01$. Family Size did not yield a reliable main effect averaged across conditions, $\beta = 0.087$, $SE = 0.061$, $t(4854) = 1.42$, $p = .155$. Neither LP, $\beta = 0.011$, $SE = 0.255$, $t(57) = 0.045$, $p = .964$, nor OS, $\beta = 0.381$, $SE = 0.350$, $t(57) = 1.09$, $p = .281$, yielded reliable main effects.
	
	The Complexity $\times$ Family Size interaction was significant, $\beta = -0.440$, $SE = 0.123$, $t(4854) = -3.58$, $p < .001$. Averaged across LP levels, the complexity effect was substantially larger for nonwords from large morphological families ($\Delta = 0.697\,\mu$V, $t = 8.02$, $p < .001$) than for those from small families ($\Delta = 0.246\,\mu$V, $t = 2.83$, $p = .005$), indicating that morphological family richness amplifies the N400 complexity effect when averaged over individual differences.
	
	The Complexity $\times$ LP interaction was also significant, $\beta = 0.250$, $SE = 0.050$, $t(4854) = 5.01$, $p < .001$. Follow-up contrasts examining the complexity effect at $-1$, $0$, and $+1$ SD of LP revealed a monotonic attenuation---the complexity effect was largest at low LP ($\Delta = 0.710\,\mu$V, $t = 9.13$, $p < .001$), attenuated at mean LP ($\Delta = 0.460\,\mu$V, $t = 7.47$, $p < .001$), and further reduced but still significant at high LP ($\Delta = 0.210\,\mu$V, $t = 2.59$, $p = .010$).
	
	The Family Size $\times$ LP interaction was also significant, $\beta = 0.121$, $SE = 0.050$, $t(4854) = 2.42$, $p = .015$, indicating that the influence of morphological family size on N400 amplitude varied as a function of LP; this interaction is qualified by the three-way interaction reported below.
	
	Critically, the three-way Complexity $\times$ Family Size $\times$ LP interaction was significant, $\beta = 0.227$, $SE = 0.100$, $t(4854) = 2.27$, $p = .023$, $\eta^2_p = .001$, indicating that LP modulated the joint influence of morphological complexity and family size on N400 amplitude. Follow-up interaction contrasts examined the Complexity $\times$ Family Size interaction separately at $-1$, $0$, and $+1$ SD of LP. At low LP, the Complexity $\times$ Family Size interaction was large and highly reliable ($\Delta = -0.667\,\mu$V, $t = -4.29$, $p < .001$)---the complexity effect was substantially larger for large-family nonwords ($\Delta = 1.044\,\mu$V, $t = 9.49$, $p < .001$) than for small-family nonwords ($\Delta = 0.376\,\mu$V, $t = 3.42$, $p = .001$), reflecting disproportionately strong semantic processing demands when both morphological complexity and family richness are high. At mean LP the interaction remained significant but attenuated ($\Delta = -0.440\,\mu$V, $t = -3.58$, $p = .0003$), with complexity effects again larger for large-family ($\Delta = 0.680\,\mu$V, $t = 7.82$, $p < .001$) than small-family nonwords ($\Delta = 0.240\,\mu$V, $t = 2.76$, $p = .006$). At high LP the interaction was no longer significant ($\Delta = -0.213\,\mu$V, $t = -1.32$, $p = .188$).
	
	To characterise the nature of this attenuation more precisely, follow-up contrasts examined the family size effect within complex nonwords across a wider range of LP, from $-2$ to $+2$ SD. At $-2$ SD of LP, large-family complex nonwords elicited substantially more negative N400 amplitudes than small-family complex nonwords ($\Delta = 0.597\,\mu$V, $t = 3.69$, $p < .001$). This advantage for large-family items decreased monotonically with increasing LP: at $-1$ SD ($\Delta = 0.367\,\mu$V, $t = 3.34$, $p = .001$), at mean LP ($\Delta = 0.133\,\mu$V, $t = 1.53$, $p = .127$), and at $+1$ SD the difference was negligible and non-significant ($\Delta = -0.102\,\mu$V, $t = -0.89$, $p = .371$). At $+2$ SD of LP the pattern reversed significantly---small-family complex nonwords now elicited more negative N400 amplitudes than large-family complex nonwords ($\Delta = -0.340\,\mu$V, $t = -1.99$, $p = .046$). This reversal indicates that for highly proficient readers, a rich morphological family facilitates rather than amplifies semantic integration demands for complex nonwords. The LP scores of participants in the present sample ranged from $-2.67$ to $3.11$, confirming that estimates at $\pm 2$ SD are grounded in observed data rather than extrapolation. Model fit indices were $R^2_{\text{marginal}} = .012$ and $R^2_{\text{conditional}} = .666$.
		
	
	
	\section{Discussion}
	
	The present study examined whether individual differences in Orthographic Sensitivity (OS) and Language Proficiency (LP) dissociably modulate morphological processing at distinct temporal stages, as indexed by the N250 and N400 components of the ERP. Three findings structure the discussion that follows. First, individual differences were largely absent from behavioural response times, with the exception of an OS advantage for nonword lexical decisions, indicating that the two dimensions differ in their contribution to general processing efficiency but that neither selectively shapes the structure of morphological effects at the behavioural level. Second, at the N250, both OS and LP modulated the complexity effect for nonword stimuli, though through qualitatively distinct mechanisms, and OS selectively modulated the base frequency effect for word stimuli, providing partial support for the predicted role of orthographic skill in early morpho-orthographic parsing. Third, at the N400, LP selectively drove the interaction between morphological complexity and family size for nonword stimuli and a crossover interaction with base frequency for word stimuli, consistent with the predicted role of LP in semantic integration. Together these findings support a partial dissociation between OS and LP in morphological processing, with the clearest dissociation emerging at the N400 and a more complex picture at the N250.
	
	\subsection{Early Morpho-orthographic Processing: The N250}
	
	\subsubsection{Word Stimuli}
	
	For word stimuli, readers higher in Orthographic Sensitivity showed larger N250 differences between high and low base frequency words, a pattern consistent with the prediction that early form-based processing is selectively shaped by the precision of sublexical orthographic representations. High base frequency words elicited more negative N250 amplitudes than low base frequency words, and this difference increased with OS, such that readers higher in orthographic sensitivity showed the largest N250 amplitudes for high frequency bases. Morphological family size also modulated the N250 as a main effect, with large-family words eliciting less negative amplitudes than small-family words, and a marginal trend toward a Family Size $\times$ OS interaction in the same direction as the base frequency finding.
	
	The opposite directions of the base frequency and family size effects  at the N250 are theoretically informative and can be accounted for by  a distinction between form-level entrenchment and semantic  reinforcement of morphological structure. Base frequency reflects how  often an orthographic form has been encountered as the same lexical  entry, across inflectional variants that share a stable meaning but do  not create new form-meaning pairings at the lexical level. Family size,  by contrast, reflects how many distinct derived words embed the base,  each constituting a separate lexical entry that pairs the base form  with a related but distinct meaning. Repeated encounter of the same  form across multiple such meaning-bearing contexts is argued to build  richer, more precise semantic representations \citep{bolgerContextVariationDefinitions2008, liLearningNovelComplex2026}; the semantic core of the base becomes more robustly  specified as it is reinforced across independent form-meaning pairings. Under this account, it is the quality of the lexical representation  associated with large family size---the precision and richness of the  semantic core built through accumulated contextual variation---that enables stronger top-down semantic feedback to the form level, facilitating orthographic identification and reducing N250 amplitude. Base frequency does not carry this semantic reinforcement: repeated encounters with a high-frequency base accumulate orthographic familiarity without necessarily building new form-meaning pairings, and its effect at the N250 therefore reflects engagement of morpho-orthographic parsing that is not necessarily supplemented by feedback from a richly specified semantic representation.
	
	The modulation of this base frequency effect by OS is consistent with the theoretical profile of this dimension. OS indexes sensitivity to sublexical orthographic structure, and readers higher in OS are more sensitive to the statistical and structural properties of sub-word units, including stem forms. For these readers, a high-frequency base strongly recruits morpho-orthographic parsing activity, producing greater N250 amplitude. The absence of a corresponding modulation by LP is consistent with the prediction that early form-based processing is specifically tied to orthographic skill rather than lexical-semantic knowledge, and supports the theoretical distinction between the two individual difference dimensions at this processing stage.
	
	\subsubsection{Nonword Stimuli}
	
	For nonword stimuli, both OS and LP modulated the N250 complexity effect, but through qualitatively distinct mechanisms that reflect the different processing profiles of the two dimensions. This constitutes a partial qualification of the predicted OS-specific modulation at the N250, but the nature of the qualification is itself theoretically informative: rather than reflecting a single shared mechanism operating at different strengths, the two dimensions produce reversals of the complexity effect through entirely distinct trajectories.
	
	For readers lower in OS, complex nonwords elicited substantially more negative N250 amplitudes than simple nonwords, reflecting a large early processing cost associated with the presence of a real suffix. As OS increased, this complexity cost attenuated monotonically. At OS = $+1$ SD the contrast was not significant ($\Delta = -0.142\,\mu$V, $t = -1.93$, $p = .053$), and at the upper extreme of the observed OS range (OS = $+1.9$ SD) the complexity effect reversed significantly, with simple nonwords eliciting more negative N250 amplitudes than complex ones ($\Delta = -0.490\,\mu$V, $t = -4.27$, $p < .001$). Examination of the estimated marginal means at each OS level revealed that this reversal was driven primarily by the trajectory of complex nonwords: both complex and simple nonwords became less negative as OS increased, reflecting a general improvement in orthographic processing efficiency, but complex nonwords decreased in negativity considerably faster than simple ones. We propose that this pattern reflects affix-led parsing: readers higher in OS are specifically sensitive to whether a suffix-position element constitutes a real affix, and this sensitivity drives increasingly efficient morpho-orthographic parsing of complex nonwords as OS increases. Simple nonwords also benefit from improved general orthographic processing efficiency, but do not benefit from affix recognition specifically, and therefore decrease in negativity more slowly. The reversal at very high OS is a consequence of the complex nonword trajectory outpacing the simple nonword trajectory — not of a rising processing cost for simple nonwords.
	
	The LP pattern was qualitatively different, and the contrast with the OS pattern constitutes a double dissociation within the nonword N250 data. As with OS, the complexity effect reversed at high LP — significantly so at LP = $+1$ SD ($\Delta = -0.141\,\mu$V, $t = -2.14$, $p = .032$) and more strongly at LP = $+2$ SD ($\Delta = -0.500\,\mu$V, $t = -5.12$, $p < .001$). However, examination of the estimated marginal means revealed that the mechanism was entirely distinct from that underlying the OS reversal. Complex nonwords were essentially flat across the LP range (EMM: $-0.70\,\mu$V at LP = $-2$ SD to $-0.82\,\mu$V at LP = $+2$ SD), indicating that LP did not modulate the processing of complex nonwords. Instead, the reversal was driven entirely by simple nonwords becoming increasingly negative as LP increased, from $+0.23\,\mu$V at LP = $-2$ SD to $-1.32\,\mu$V at LP = $+2$ SD. We propose that this pattern reflects stem-led parsing: readers higher in LP have richer and more precise representations of stems as free morphemes with discrete lexical entries, and recognition of the stem in a nonword triggers an attempt to assign morphological structure to the whole string regardless of what follows. For complex nonwords, this parse succeeds --- a real stem is followed by a real suffix, yielding a valid morphological analysis --- and successful parsing produces neither a cost nor a benefit at the N250, keeping amplitude stable across LP levels. For simple nonwords, the parse fails --- the pseudoaffix cannot be integrated into a valid morphological structure --- and this failed parse produces a processing cost that grows with LP, driving the simple nonword trajectory increasingly negative.
	
	The contrast between the two reversal mechanisms is theoretically precise. The OS reversal is driven by the complex nonword trajectory --- affix recognition produces increasingly efficient parsing of complex forms, while simple nonwords are not specifically affected. The LP reversal is driven by the simple nonword trajectory --- stem recognition triggers parsing attempts that succeed for complex nonwords without modulating their amplitude, but fail for simple nonwords at a cost that grows with LP. This dissociation maps directly onto the theoretical profiles of the two dimensions: OS indexes sensitivity to the sublexical orthographic units that affixes represent, and its parsing signature is carried by the condition in which a real affix is present; LP indexes depth of lexical-semantic knowledge in which stems are embedded as free morphemes, and its parsing signature is carried by the condition in which stem recognition triggers a parse that cannot be resolved. Together, these findings suggest that early morpho-orthographic processing of novel forms is sensitive to both affix-based and stem-based knowledge, with each dimension leaving a distinct fingerprint in the N250 through qualitatively different parsing mechanisms.
	
	subsection{Semantic Integration and the N400}
	
	The N400 findings provide the clearest support for the predicted dissociation between Orthographic Sensitivity and Language Proficiency. Across both word and nonword stimuli, LP selectively modulated N400 amplitude in interaction with morphological variables, while OS appeared only as a main effect reflecting general differences in the efficiency of semantic access independently of morphological structure. This pattern is consistent with the theoretical prediction that LP, as an index of the depth and precision of lexical-semantic representations, selectively shapes the semantic consequences of morphological processing, while OS, as an index of sublexical orthographic sensitivity, does not specifically modulate semantic integration.
	
	\subsubsection{Word Stimuli}
	
	For word stimuli, LP modulated the N400 effect of base frequency in a manner that crossed over across the LP continuum. Before unpacking this finding, it is worth stating plainly what it reveals: the relationship between base frequency and the ease of semantic access during word recognition is not fixed, but depends on where a reader falls on the LP continuum. For readers lower in Language Proficiency, words with higher base frequency were associated with greater N400 negativity than words with lower base frequency --- a pattern suggesting that for these readers, base frequency acts as a familiarity signal that gates the degree to which semantic access is attempted. High frequency bases, being more orthographically familiar, trigger more active semantic access attempts, while low frequency bases, being less familiar, engage semantic access only partially. For readers higher in Language Proficiency, this pattern reversed: high frequency bases were associated with less negative N400 amplitudes than low frequency bases, consistent with the more familiar orthographic form supporting more efficient mapping onto its semantic representation.
	
	This crossover is consistent with the Lexical Quality Hypothesis \citep{Perfetti2007}, under which higher quality lexical representations involve more precise and fully specified form-to-meaning bindings. For low LP readers, whose representations are less precise, base frequency operates as a crude familiarity heuristic that determines the depth of semantic engagement. For high LP readers, whose representations are more fully specified, base frequency influences the efficiency of form-to-meaning mapping rather than gating whether mapping is attempted at all. We note, however, that the linear interaction term in the model does not distinguish between these two accounts --- the data are consistent with a continuous modulation of the contribution of base frequency to semantic access across the LP continuum, and the nature of that modulation cannot be fully characterised from the present data alone. Accuracy data would be needed to adjudicate between an account in which low LP readers partially disengage semantic access for low frequency bases and one in which they engage it but less efficiently.
	
	\subsubsection{Nonword Stimuli}
	
	The nonword N400 finding is the centrepiece result of the present study, and its interpretation is most readily approached by building from the simpler two-way interaction to the full three-way. At the simplest level, the Complexity $\times$ Family Size interaction, averaged over LP, indicates that the N400 complexity effect --- the greater negativity for complex than simple nonwords --- was substantially larger for nonwords built on stems with large morphological families than for those built on stems with small families. This indicates that morphological family richness amplifies the semantic processing demands associated with encountering a morphologically complex novel form: a complex nonword built on a stem with many morphological relatives activates a richer semantic network, producing greater N400 negativity than one built on a stem with fewer relatives.
	
	The critical question is what LP does to this two-way interaction --- and here the finding can be stated plainly before the statistics are unpacked: for readers lower in Language Proficiency, a rich morphological family amplifies semantic processing demands for complex nonwords; for readers higher in Language Proficiency, that same richness becomes a resource that facilitates processing (see Figure~\ref{fig:nw_n400_threeway}). The Complexity $\times$ Family Size interaction was large and significant at low LP, attenuated at mean LP, and absent at high LP, with the family size effect within complex nonwords reversing significantly at LP $= +2$ SD. Low LP readers showed disproportionately strong N400 negativity for large-family complex nonwords relative to small-family ones, reflecting greater semantic processing demands when both morphological complexity and family richness are high. For these readers, the rich semantic network activated by a large-family stem creates demands that cannot be efficiently resolved --- more semantic content is recruited in an attempt to integrate a form that has no lexical entry, producing greater negativity. High LP readers, by contrast, showed the reverse pattern at LP $= +2$ SD: large-family complex nonwords elicited less negative N400 amplitudes than small-family ones, indicating that for these readers, morphological family richness facilitates rather than amplifies semantic processing demands.
	
	This reversal is theoretically interpretable within the Lexical Quality Hypothesis. High LP readers have more precise and fully specified lexical-semantic representations, including richer representations of stems and their morphological relatives. When such a reader encounters a complex nonword built on a large-family stem, the richness of the morphological neighbourhood provides a scaffold for processing --- the activated semantic network is well-organised and precisely specified, allowing the reader to efficiently evaluate the novel form against known morphological structure. For a reader with less precise representations, the same activated network constitutes a processing burden rather than a scaffold, because the semantic content cannot be efficiently organised or resolved against the novel form.
	
	This account is also consistent with the distinction between base frequency and family size developed in the context of the N250 findings. Family size reflects the richness of the semantic core built through repeated encounter of a stable form with a stable meaning across multiple distinct form-meaning pairings \citep{bolgerContextVariationDefinitions2008, liLearningNovelComplex2026}. It is precisely this semantic richness that makes family size a double-edged variable at the N400: for readers whose representations are insufficiently precise to leverage it, semantic richness amplifies processing demands; for readers whose representations are sufficiently precise, it provides a processing advantage. Base frequency, which accumulates orthographic familiarity without building new form-meaning pairings, does not carry this semantic richness, and its role at the N400 is correspondingly different --- modulating the efficiency of form-to-meaning mapping rather than the depth of semantic network activation. Together, the word and nonword N400 findings converge on a coherent picture: LP indexes the precision of lexical-semantic representations, and it is this precision that determines whether morphological information --- whether the orthographic entrenchment of a familiar base or the semantic richness of a large morphological family --- supports or burdens semantic access during word recognition.
	
	\subsection{Figures}
	Frontiers requires figures to be submitted individually, in the same order as they are referred to in the manuscript. Figures will then be automatically embedded at the bottom of the submitted manuscript. Kindly ensure that each table and figure is mentioned in the text and in numerical order. Figures must be of sufficient resolution for publication \href{https://www.frontiersin.org/about/author-guidelines#ImageSizeRequirements}{see here for examples and minimum requirements}. Figures which are not according to the guidelines will cause substantial delay during the production process. Please see \href{https://www.frontiersin.org/about/author-guidelines#FigureRequirementsStyleGuidelines}{here} for full figure guidelines. Cite figures with subfigures as figure \ref{fig:Subfigure 1} and \ref{fig:Subfigure 2}.
	
	
	\subsection{Tables}
	Tables should be inserted at the end of the manuscript. Please build your table directly in LaTeX.Tables provided as jpeg/tiff files will not be accepted. Please note that very large tables (covering several pages) cannot be included in the final PDF for reasons of space. These tables will be published as \href{http://home.frontiersin.org/about/author-guidelines#SupplementaryMaterial}{Supplementary Material} on the online article page at the time of acceptance. The author will be notified during the typesetting of the final article if this is the case. 
	
	
	\section{Additional Requirements}
	
	For additional requirements for specific article types and further information please refer to \href{http://www.frontiersin.org/about/AuthorGuidelines#AdditionalRequirements}{Author Guidelines}.
	
	\section*{Conflict of Interest Statement}
	The authors declare that the research was conducted in the absence of any commercial or financial relationships that could be construed as a potential conflict of interest.
	
	\section*{Author Contributions}
	JM: Conceptualization, Methodology, Software, Validation, Formal analysis, Resources, Data curation, Visualization, Writing – original draft, Writing – review \& editing, Supervision.
	EK: Methodology, Software, Investigation, Data curation, Writing – original draft, Writing – review \& editing, Supervision.
	FG: Methodology, Software, Investigation, Data curation.
	JP: Investigation, Data curation, Writing – review \& editing
	TR: Investigation, Data curation, Writing – review \& editing
	EM: Investigation, Data curation, Writing – review \& editing
	CR-M: Investigation, Data curation, Writing – review \& editing
	NA: Investigation, Data curation, Writing – review \& editing
	
	All authors contributed to the article and approved the submitted version.
	
	\section*{Funding}
	Research reported in this manuscript was supported by the Rhode Island IDeA Network of Biomedical Research Excellence, funded by the National Institute of General Medical Sciences of the National Institutes of Health under grant number P20GM103430. Additional support for student researchers (Joemari Pulido, Tess Rooney, Natalia Alzate, Ethan Moore, and Crismar Ramos-Marte) and for data collection was provided through the Center for Engaged Learning at Providence College and a series of Hampshire College Faculty Development Grants awarded to Dr. Morris during her tenure at that institution. The funders had no role in study design, data collection and analysis, decision to publish, or preparation of the manuscript.
	
	
	\section*{Acknowledgments}
	We thank Sophie Fulghum, Dylan Palumbo, Maeve Garrity, Aleksandra Kurylowicz, Sarah Debian, Aleksandra Saneff, Kenny Diaz, and Kiley Vallee for their assistance with participant recruitment and data collection.
	
	\section*{Supplemental Data}
	\href{http://home.frontiersin.org/about/author-guidelines#SupplementaryMaterial}{Supplementary Material} should be uploaded separately on submission, if there are Supplementary Figures, please include the caption in the same file as the figure. LaTeX Supplementary Material templates can be found in the Frontiers LaTeX folder.
	
	\section*{Data Availability Statement}
	The datasets [GENERATED/ANALYZED] for this study can be found in the [NAME OF REPOSITORY] [LINK].
	% Please see the availability of data guidelines for more information, at https://www.frontiersin.org/about/author-guidelines#AvailabilityofData
	
	\bibliographystyle{Frontiers-Harvard} % Many Frontiers journals use the Harvard referencing system (Author-date), to find the style and resources for the journal you are submitting to—https://zendesk.frontiersin.org/hc/en-us/articles/360017860337-Frontiers-Reference-Styles-by-Journal. For Humanities and Social Sciences articles please include page numbers in the in-text citations 
	%\bibliographystyle{Frontiers-Vancouver} % Many Frontiers journals use the numbered referencing system, to find the style and resources for the journal you are submitting to—https://zendesk.frontiersin.org/hc/en-us/articles/360017860337-Frontiers-Reference-Styles-by-Journal
	\bibliography{M21_LDT.bib}
	
	
	\end{document}